\documentclass[11pt,a4paper]{article}
\usepackage[T1]{fontenc}
\usepackage[utf8]{inputenc}
\usepackage[margin=25mm]{geometry}
\usepackage{lmodern}
\usepackage{microtype}
\usepackage{booktabs}
\usepackage{amsmath}
\usepackage[hidelinks]{hyperref}
\usepackage[numbers,sort&compress]{natbib}
\setlength{\parindent}{0pt}
\setlength{\parskip}{0.55em}
\hyphenpenalty=10000
\exhyphenpenalty=10000
\setlength{\emergencystretch}{2em}
\newcommand{\paperauthors}{Authors to be confirmed}
\newcommand{\paperrepository}{\href{https://github.com/yunusserhat/firerisk}{github.com/yunusserhat/firerisk}}
\hypersetup{pdftitle={SigLIP2 for aerial fire risk classification},pdfauthor={\paperauthors}}
\title{SigLIP2 for aerial fire risk classification}
\author{\paperauthors}
\date{}

\begin{document}
\maketitle

\begin{abstract}
We examine the transfer of a pretrained SigLIP2 image encoder to seven class
fire risk classification from aerial imagery. We introduce a reproducible
partition of the public FireRisk training mirror and an implementation that
records data provenance, preprocessing and model selection. Two initial runs
compare a frozen encoder probe with full model adaptation. On the validation
partition, full adaptation reaches 63.05\% accuracy and 58.94\% macro F1,
compared with 55.95\% and 50.19\% for the probe. Both runs use one training
seed and select their checkpoint on the same validation partition. These
development results support further evaluation of SigLIP2 but do not establish
performance on an independent test set or unseen regions. The accompanying
code provides a common framework for repeated experiments and comparisons
with additional visual encoders.
\end{abstract}

\section{Introduction}

Aerial images describe vegetation, buildings and surface cover that may help
distinguish wildfire hazard categories. FireRisk links such imagery to labels
derived from the Wildfire Hazard Potential map~\citep{shen2023firerisk}. The
task is to recover a hazard category assigned to an image tile. It does not
directly measure the probability of a future fire at that location.

SigLIP2 learns visual representations using image and text supervision together
with additional representation learning objectives~\citep{tschannen2025siglip2}.
This motivates an empirical question about whether its image encoder can
adapt to the different visual content of aerial imagery. We study a small
initial comparison and release its protocol and implementation for further
experiments.

\section{Data and protocol}

The original FireRisk study describes 91,872 images, with 70,331 training
examples and a separate validation partition~\citep{shen2023firerisk}. Its RGB
imagery was obtained from the National Agriculture Imagery Program during
2019 and 2020 and cropped to $320\times320$ pixels. Labels come from the
2020 Wildfire Hazard Potential map. We use the
\href{https://huggingface.co/datasets/blanchon/FireRisk}{\texttt{blanchon/FireRisk}}
mirror, which exposes only the 70,331 training examples. The mirror therefore
does not provide the published evaluation partition.

We create a new split with 49,231 training images, 10,552 validation images
and 10,548 reserved test images. Allocation is stratified by class, using
split seed 2026. The seven categories are very low, low, moderate, high, very
high, nonburnable and water. The split preserves the class imbalance. For
example, the validation partition contains 3,264 very low examples and 490
very high examples.

Before allocation, the preparation code hashes decoded RGB pixels and image
dimensions using SHA256 after correcting image orientation. It groups exact
duplicates and rejects conflicting labels within a group. The audit found
70,331 unique pixel groups and no repeated images in this revision. This
check concerns exact identity. It cannot exclude overlapping tiles or visual
similarity between nearby locations. The mirror contains image and label
fields without coordinates or acquisition dates, so this experiment cannot
establish geographic or temporal independence.

The data revision, model revision and split hash are recorded with each run.
Validation macro F1 selects the checkpoint. The reserved test partition has
not been evaluated in this study. As a result, our scores are development
measurements under a new protocol and cannot be compared directly with the
original FireRisk benchmark.

\section{Model and training}

We use the SigLIP2 base image encoder with $16\times16$ patches at
$224\times224$ input resolution, distributed through
\href{https://huggingface.co/timm/vit_base_patch16_siglip_224.v2_webli}{timm}.
The pretrained attention pooling is retained. A trainable LayerNorm and a
seven output linear layer operate on its 768 dimensional representation.
The resulting classifier has 92,891,143 parameters. In the frozen encoder
probe, only the head changes, leaving 6,919 trainable parameters. The encoder
remains in evaluation mode during training. Full adaptation updates all
parameters and uses head dropout with probability 0.1.

Training uses random crops covering 80\% to 100\% of the image area, with
aspect ratios from 0.9 to 1.1. Horizontal and vertical flips have probability
0.5, and rotations use multiples of 90 degrees. Evaluation uses deterministic
bicubic resizing and the pretrained center crop policy with crop proportion
0.9. Each RGB channel is normalized with mean 0.5 and standard deviation 0.5.

Both runs use AdamW with weight decay 0.05, label smoothing 0.05 and gradient
clipping at norm 1. Learning rates follow a cosine schedule after warmup over
10\% of the configured optimization steps. The probe uses a head learning
rate of $10^{-3}$ and at most 15 epochs. Full adaptation uses $10^{-5}$ for
the encoder and $3\times10^{-4}$ for the head, with at most 30 epochs.
No class weighting or mixup is applied. The effective batch size is 64 in
both cases. Full adaptation accumulates gradients over two batches of 32
images and uses gradient checkpointing. Computation uses bfloat16 on one
NVIDIA RTX 5090 GPU. Training seed 42 is shared by the two runs.

Early stopping follows seven epochs without a validation macro F1 improvement
greater than 0.0001. Both runs stop after 14 epochs and select epoch 7. The
different learning rates, dropout and schedule lengths make this a comparison
of two initial recipes. They prevent attribution of the observed difference
to encoder adaptation alone.

\section{Preliminary results}

Table~\ref{tab:results} reports uncalibrated predictions from the selected
checkpoints. Full adaptation improves validation accuracy by 7.10 percentage
points and macro F1 by 8.75 percentage points relative to the probe. Macro F1
assigns equal weight to the seven class F1 scores, making the minority classes
visible alongside accuracy.

\begin{table}[!htb]
\centering
\small
\refstepcounter{table}\label{tab:results}
\begin{tabular}{lrrrrr}
\toprule
Recipe & Trainable & Accuracy & Macro F1 & Time & Memory \\
       & parameters & (\%) & (\%) & (min) & (GiB) \\
\midrule
Frozen encoder probe & 6,919 & 55.95 & 50.19 & 12.1 & 0.74 \\
Full adaptation & 92,891,143 & 63.05 & 58.94 & 19.0 & 1.96 \\
\bottomrule
\end{tabular}
\medskip
\parbox{\textwidth}{Table~\thetable. Validation results from one training
seed. Time covers the training
loop, validation and checkpoint writing, excluding data preparation and
model loading. Memory is peak PyTorch allocated GPU memory and does not
represent total device or system memory.}
\end{table}

The strongest class F1 scores after full adaptation are water at 87.33\%
and nonburnable at 85.12\%. Very high reaches 39.90\% F1 and 33.47\% recall,
compared with 21.29\% F1 and 14.49\% recall for the probe. Although this
minority category improves, its recall remains low. Full adaptation correctly
identifies 164 of its 490 validation examples. This limits the usefulness of
the current classifier when missing high hazard areas has a substantial cost.

The runs executed concurrently on separate RTX 5090 GPUs. Their timings are
contextual measurements and do not provide a controlled speed comparison.
Times exclude environment installation, downloads and the initial pixel
audit. No repeated training seeds are available, so the variation caused by
training randomness remains unknown.

\section{Reproducibility and limitations}

The implementation keeps original Parquet shards and can build an indexed
archive of unchanged image bytes for shuffled access. It avoids extracting
each image to a separate file and records the selected artifact location.
Each run saves its configuration, environment, preprocessing, data manifest,
training history and selected weights. Immutable dataset and model revisions
and the split hash are listed in the accompanying manuscript documentation.
The initial run artifacts do not record a source commit and mark the
development tree as modified.
The repository also includes presets for LoRA adaptation, DINOv2, ConvNeXt
and ResNet50. Those presets have no completed comparative results here.

Several steps remain before a broader empirical claim is warranted. The
recipes should receive declared validation search budgets and be repeated
with independent training seeds. The fixed recipes can then be evaluated on
the reserved test partition. A stronger study of transfer to new locations
requires geographic metadata and a separate spatial holdout. Exact pixel
hashes do not resolve geographic correlation or possible overlap with the
pretraining data.

Full SigLIP2 adaptation merits further study on this mirror. Errors within
the risk categories, particularly very high hazard, remain a central concern.

\paragraph{Code availability}
\paperrepository. The code is distributed under the GNU General Public
License version 3. Upstream imagery and pretrained weights retain their
respective terms.

\bibliographystyle{plainnat}
\bibliography{references}

\end{document}
